\originalchapter{离开主序、核壳分离与巨星演化}
本章支撑8.901第10--12及16次课。恒星演化的关键不是在H--R图上背诵箭头，而是追踪哪一层还有燃料、哪一层供压、热从哪里进入和离开。前置为结构方程、核产热和简并压；本章的图示是定性阶段关系，数值轨迹仅在指定作业模型里给出。

\originalsection{核心氢耗尽改变了热源的位置}
主序核心逐渐积累氦，$\mu$增大，在相同压强、温度下气体密度较高。最亮区域的氢耗尽后，核心没有相同的核热源，向外损失能量会收缩并升温；其外仍富氢的壳可达到燃烧温度，热源从中心迁到壳层。新的氦灰不断加入核心，使核质量增长。

外包层膨胀不应简单归结为“核心更热所以整星更热”。核心收缩、壳层光度、包层压强和熵同时调整；包层半径增大后，$T_{\rm eff}=[L/(4\pi R^2\sigma_{\rm SB})]^{1/4}$可降低，即使总光度提高。究竟沿H--R图哪条曲线走，需要解质量分层的能量方程，不能仅凭全星维里关系唯一预测半径。

全对流星的混合可以把更多氢送入核心；有辐射区的星则有明显核壳组成分层。这解释为何“可燃质量比例”不是全质量的固定百分数。PS5.1先给各分支的燃料参数，PS5.2再对明确给定的核心质量--光度--半径经验关系实际积分；后者的适用区间也需检查。

\originalsection{Schönberg--Chandrasekhar极限是匹配极限}
氢耗尽后，一个无核产热、非简并、近等温核心可暂以理想气体供压。它要和包层在同一半径匹配压强、温度与包围质量；两侧$\mu$不同，同压、同温使$\rho_c/\rho_e=\mu_c/\mu_e$。允许的核解不是任意大的等温球都能接到同一个包层。

引入同源变量$U=\dd\ln m/\dd\ln r=4\pi r^3\rho/m$、$V=-\dd\ln P/\dd\ln r=Gm\rho/(rP)$。核、壳的静水平衡与质量方程可在$U,V$平面追踪。界面$\rho$跳变时$U,V$也按$\mu_c/\mu_e$相应变化；可匹配的包层轨迹只覆盖有限区域。增加核质量使核轨迹最终离开可接区域，产生最大核质量比例。这是边界值问题的极限，不是把某层温度简单设成零。

尺度可从等温核的维里预算看出。设界面压强$P_s$、核温$T_c$，其质量尺度只能由$G,P_s,c_T^2=k_BT_c/(\mu_cm_u)$组成：
\[
 M_c\sim\frac{c_T^4}{G^{3/2}P_s^{1/2}}.
\]
相同界面状态的气压包层有类似质量尺度、但含$\mu_e$，所以最大核比例的方向约为$q_{\max}\propto(\mu_e/\mu_c)^2$。无量纲系数要实际求核--包层匹配，以上量纲推导没有计算它。

已查\href{https://articles.adsabs.harvard.edu/pdf/1942ApJ....96..161S}{Schönberg与Chandrasekhar的1942原论文}：其指定的等温核、辐射包层、核／包层分子量比2的模型，最大核质量比例约0.101。此为该历史模型的匹配结果，不能称作所有恒星的普遍10\%上限；不同包层、多方近似、辐射压或简并会改它。本卷没有把未重新求解的匹配系数声称为独立数值验证。

超过这个非简并等温匹配极限，核心需改变状态并继续收缩，直到新的热源或压强支持出现；它不等于全星马上自由落体，也不等于电子简并的Chandrasekhar总质量上限。

\originalsection{简并氦核与非简并氦点火走不同路线}
较低质量恒星的氦核可在点火前变简并，压强主要由密度控制。三氦率对温度陡峭，初升温无法迅速靠气压膨胀调节，因而局部热失控，最终升温、扩张解除简并并恢复热调节。称“氦闪”主要说核心事件，不能从局部产热率直接推表面出现同样短、同样强的光闪。

较高质量核心在点火前较热、非简并，燃烧和膨胀反馈可较平稳。氦核心烧完留下碳氧核心后，又有氦和氢壳燃烧。薄壳的几何和温度敏感性可产生热脉冲；质量损失将外包层送入星际介质，留下白矮星的条件取决于核质量、组成和演化史，不由一个H--R点唯一决定。

沿定性阶段链，可写核心H燃烧$\to$惰性He核与H壳$\to$核心He燃烧$\to$惰性CO核及双壳；大质量分支还能进入更重燃烧阶段。每一步都应同时检查燃料、EOS与时间尺度，避免将观测分类如“红巨星”误当一种唯一内部核反应。

\originalsection{年龄信息来自转折和组成，而非图上一条直线}
单星的$L,T_{\rm eff}$可能对应多个阶段；双星相互作用还会改变质量与可用燃料。近同龄星团的主序转折告诉我们哪一初始质量已耗尽核心氢，结合$L(M)$和可燃质量可估年龄。金属丰度改变不透明度与谱，混合改变核心大小，均影响时间；因此“转折温度”与年龄之间需一族指定物理输入的模型。

H--R轨迹若按给定关系实际积分，必须标明固定质量、经验公式范围与时间参数。PS5.2输出的四种图遵循原题低核质量起点，有很长的早段；没有把这种教学曲线假装成已经运行专业恒星演化代码。
